\section{Identification algorithm: complete utility, termination and soundness proofs}
\label{sec:annex-p16-identification-algorithm}

\begin{annexitem}{annex:p16.id-positive-support}{thm:p16-id-positive-support}{Positivity in the current and further-intervened models}
\begin{proof}
Fix the outside parameters of a current model. Its full law \(q_R\)
is positive at every observed assignment. The latent-input formula
expresses its mass as a sum of prior latent weights times all
retained equation indicators. Fixing an additional observed set
\(I\subseteq R\) removes its indicators. At a full assignment
whose \(I\) coordinates equal the proposed action values, this
can only increase each nonnegative summand. Thus the intervention
mass on the remaining coordinates is at least \(q_R(v)>0\).
Every partial observation disjoint from the intervention has a full
extension, since all finite value sets are nonempty, and its marginal
probability is therefore positive as well.

An ancestral marginal is a finite sum with at least one positive
summand at every retained assignment and outside parameter. For a
positive kernel, every prefix marginal is positive, and every ratio
of consecutive prefixes is positive. Their district product is
therefore positive. By the proved district recovery and restriction
theorems it equals the normalized kernel \(Q[D]\) of the restricted
model, with the newly fixed arguments retained. These facts prove
the positivity claims for all operations and all do-calculus
conditionals below, including empty observation sets whose mass is one.
\end{proof}
\end{annexitem}

\begin{annexitem}{annex:p16.id-do-licenses}{thm:p16-id-do-licenses}{Exact path arguments behind the source utility derivations}
\begin{proof}
Work in the expanded DAG of the current model, with independent
latent roots and observed vertices \(R\). A directed path that
starts at an observed vertex cannot enter a latent root, since roots
have no incoming edges. Directed ancestry between observed vertices
is consequently exactly ancestry in the directed part of \(\Gamma\),
also after deleting incoming or outgoing observed arrows.

Each local rule is an original-model rule with the current fixed set
\(F=N\setminus R\) adjoined to its baseline action set. In the
original expanded graph after fixing \(F\), those vertices have
no incoming edges. Conditioning on their fixed values blocks any
path passing through them as internal noncolliders, and they cannot
be colliders. Nor can they activate any other collider, because no
other vertex is their ancestor. Thus active paths between retained
vertices, with the prescribed local conditioning, are exactly those
in the restricted expanded graph. Directed paths from a nonfixed
vertex to the local conditioning set cannot pass through \(F\)
either, so the Rule 3 ancestor exception is preserved. The local
conditional law \(P^R_x(y\mid w)\) is the original law
\(P_{v_F,x}(y\mid w)\). This proves the legitimacy of lifting
the local do-calculus steps back to the original model.

We use one elementary path fact. Every vertex on an active path is
an ancestor of an endpoint or of the conditioning set. A collider
is such an ancestor by activity. At a noncollider, follow an outgoing
edge along the path until an endpoint or a collider is reached.
This gives the requisite directed path to one of those vertices.
This argument is the active-path part of the already proved moral
separation bridge, and applies to the expanded DAG itself.

\paragraph{Deletion of actions at nonancestors.}
Put \(K=\widehat\Gamma_{\overline X}\). Under the first
hypothesis, no member of \(Z\) is an ancestor of \(Y\cup W\)
in \(K\), so the exact exception set for Rule 3 is all of
\(Z\). In \(K_{\overline Z}\), every proposed action vertex
has no incoming edge. A path starting there must begin with an
outgoing arrow. If it reaches \(Y\) while all arrows continue
in that direction, it contradicts the nonancestor hypothesis.
Otherwise its first reversal is a collider descended from that
starting vertex. The collider cannot be an ancestor of \(W\),
which would make the starting vertex an ancestor of \(W\).
It cannot be an ancestor of \(X\), because those conditioned
vertices have no incoming edges in \(K\); a collider cannot
itself be one of those roots. Thus the collider is inactive given
\(X\cup W\). Every such path is blocked, giving exactly the
Rule 3 graph condition. The support lemma supplies both conditional
probabilities, so the action may be inserted or deleted.

\paragraph{Insertion of earlier nonancestor observations.}
Use \(A_D,B_i,D_i,E_i\) from the statement. Suppose there is
an active path from \(E_i\) to \(i\) given \(B_i\cup D_i\)
in \(\widehat\Gamma_{\overline{B_i}}\). By the path fact,
every observed vertex on it is at or before \(i\): all observed
endpoints and conditioned vertices are at or before \(i\), and
a later observed vertex cannot be their directed ancestor.
An internal \(B_i\) vertex cannot occur, since it is now a
root and is conditioned; it is a noncollider and blocks the path.
Any internal vertex in \(D\) therefore lies in \(D_i\) and
must be a collider because it is conditioned. All other observed
vertices on the path belong to \(E_i\).

Consider its first entry into an observed vertex of \(D\).
The preceding observed vertex is in \(E_i\). The connection
between the two observed vertices is either a directed observed
edge or a fork through one latent root; two latent roots cannot
be adjacent. A directed edge from \(E_i\) into \(D\) is
impossible, because it would make that \(E_i\) vertex an
ancestor of \(D\). A directed edge from \(D\) to \(E_i\)
has a tail at the \(D\) vertex. If this vertex is internal,
it is a conditioned noncollider and blocks the path; if it is
the endpoint \(i\), the edge points to an earlier observed
vertex, contrary to the order. Finally, a shared latent root
would create a bidirected edge joining \(E_i\) to \(D\),
contradicting maximality of \(D\). Every possibility is excluded.
This proves the displayed separation and the Rule 1 license.

\paragraph{Exchange of earlier actions and observations.}
Let \(B\subseteq L_i\setminus D\), and suppose an active
path joins \(B\) to \(i\) given \(L_i\setminus B\) in
\(\widehat\Gamma_{\underline B}\). Again all observed vertices
on it are at or before \(i\). An internal vertex of \(B\)
would have two incoming arrows because its outgoing arrows were
deleted. It is an unconditioned collider with no descendants and
is therefore inactive. All remaining internal observed vertices
are conditioned and must be colliders.

There can be no directed observed edge on this path. Its tail
cannot be an internal conditioned vertex, which would then be
a blocked noncollider. Its tail cannot be the endpoint in \(B\),
whose outgoing arrows are deleted. Nor can its tail be \(i\),
since its other endpoint would be earlier in the order. Thus
consecutive observed vertices along the path can only share
latent roots. They give a bidirected path in \(\Gamma\) from
\(B\) to \(i\), contrary to \(B\cap D=\varnothing\).
This proves the Rule 2 license. All empty action/observation cases
are included by absence of endpoint paths. The support conditions
in every application follow from \cref{thm:p16-id-positive-support}.
\end{proof}
\end{annexitem}

\begin{annexitem}{annex:p16.id-do-utility}{thm:p16-id-do-utility}{The source Lemmas 4--8 route with every graphical step supplied}
\begin{proof}
The following derivations retain the do-calculus route of the source
\cite[Lemmas~4--8, p.~1224]{ShpitserPearl2006JointIdentification}.
All statements are within a current restricted model, with its
outside parameters held fixed. Its positive law and the support
lemma make every conditional used here defined.

\paragraph{Lemma 4: district factors and the product.}
Fix a maximal district \(D\), and put
\(A_D=\operatorname{an}_{\Gamma}(D)\setminus D\), with
\(a_D=v_{A_D}\).
Vertices in \(R\setminus(D\cup A_D)\) are nonancestors of
\(D\). The first license in \cref{thm:p16-id-do-licenses}
therefore removes their interventions by Rule 3, giving
\[
 \mathcal Q_R[D]=P^R_{v_{R\setminus D}}(v_D)=P^R_{a_D}(v_D).
\]
The ordinary chain rule, in the order restricted to \(D\), gives
\[
 P^R_{a_D}(v_D)=\prod_{i\in D}P^R_{a_D}(v_i\mid v_{D_i}).
\]
For each factor split \(A_D\) into the earlier vertices \(B_i\)
and its vertices later than \(i\). The latter cannot be ancestors
of \(i\) or of \(D_i\), so Rule 3 and the first license delete
their actions. The second license then inserts the observations
\(E_i\) by Rule 1. Finally apply the third license with
\(B=B_i\), noting \(L_i\setminus B_i=D_i\cup E_i\), to
exchange those earlier actions with observations by Rule 2. Thus
\[
\begin{aligned}
 P^R_{a_D}(v_i\mid v_{D_i})
 &=P^R_{b_i}(v_i\mid v_{D_i})\\
 &=P^R_{b_i}(v_i\mid v_{D_i},v_{E_i})\\
 &=q_R(v_i\mid v_{L_i}).
\end{aligned}
\]
This explicitly justifies the earlier nonancestor observations implicit
in the source's chain expression. Multiplying gives
\(\mathcal Q_R[D]=\operatorname{Ext}^R_D(q_R)\).
Multiplying this identity over all maximal districts covers every
ordered observed vertex once. The ordinary full chain rule therefore
gives \(\prod_D\mathcal Q_R[D]=q_R\).

Apply this same derivation to the model after intervention \(X=x\).
Its remaining observed graph is \(\Gamma[S]\), where
\(S=R\setminus X\), and its observed law is positive by the
support lemma. Its district factors, after combining the baseline
action with the complementary district actions, are precisely
\(\mathcal Q_R[D]\) for \(D\in\mathcal D_S\). Summing the
product over \(S\setminus Y\) proves the first identity of
the statement for arbitrary action and response sets. This use of
the post-action law is a symbolic derivation of the target; it does
not assume that law was already recovered from observations.

\paragraph{Lemma 5: restriction to original ancestors.}
Let \(A=\operatorname{an}_{\Gamma}(Y)\), and let
\(W=R\setminus A\). These vertices are nonancestors of all
of \(A\). Rule 3 with the first license gives
\(P^R_w(v_A)=q_R(v_A)=\Sigma_W q_R\). This is the distribution
of the restricted model on \(A\). Write \(X_A=X\cap A\).
Members of \(X\setminus A\) are nonancestors of \(Y\), even
after fixing \(X_A\). Rule 3 first removes their actions and
then inserts actions on all of \(W\), giving
\[
 P^R_x(y)=P^R_{x_{X_A}}(y)
        =P^R_{x_{X_A},w}(y)=P^A_{x_{X_A}}(y).
\]
The first equality uses the original values where needed; the second
permits arbitrary values on the other \(W\) coordinates. The
restricted distribution and target are independent of those values.

\paragraph{Lemma 6: additional actions after fixing X.}
Put \(S=R\setminus X\) and
\(W=S\setminus\operatorname{an}_{\Gamma[S]}(Y)\).
After incoming arrows to \(X\) are deleted, a directed path
starting outside \(X\) cannot pass through \(X\). Hence
these \(W\) vertices have no directed path to \(Y\) in that
post-action graph. The first license with empty retained observations
gives Rule 3's exact premise and proves
\(P^R_x(y)=P^R_{x,w}(y)\) for every \(w\).

\paragraph{Lemma 7: the maximal-district base.}
If \(S=R\setminus X\) is already a maximal district of
\(\Gamma\), its complementary intervention kernel is
\(\mathcal Q_R[S]\). The Lemma 4 derivation above identifies
it as \(\operatorname{Ext}^R_S(q_R)\). Marginalizing
\(S\setminus Y\) proves the fourth identity.

\paragraph{Lemma 8: recursion within a larger district.}
Suppose \(S\subsetneq D\in\mathcal D_R\). The Lemma 4
derivation identifies \(q_D=\mathcal Q_R[D]\). It is the
distribution of the restricted model after fixing \(R\setminus D\),
with those assignments retained. Since \(S=R\setminus X\)
is contained in \(D\), one has
\(R\setminus D\subseteq X\) and
\(X=(R\setminus D)\mathbin{\dot\cup}(X\cap D)\).
The union symbol with a dot denotes this disjoint union. Replacing
these disjoint equation sets successively is exactly their joint
replacement. Consequently the original target is the target under
\(X\cap D\) in this restricted model, at the original values
on \(R\setminus D\). This proves the last identity. No general
latent projection, independent pairwise replacement, or unproved
identification conclusion was used in these derivations.
\end{proof}
\end{annexitem}

\begin{annexitem}{annex:p16.id-invariants}{thm:p16-id-invariants}{Valid kernels, exhaustive branches and the retained failure facts}
\begin{proof}
Initially \(R=N\), \(q_R=p=Q[N]\), and all call conditions
are assumptions. Induct along recursive edges. At branch 2,
\(A\) contains \(Y\) and is ancestral within \(R\).
The foundation's ancestral identity gives
\(\Sigma_{R\setminus A}q_R=Q[A]\), and the support lemma
makes it strictly positive. The restriction theorem gives graph
\(G[A]\) with the original root incidences intersected with \(A\).
Its fixed set is \(N\setminus A\); the newly fixed coordinates
\(R\setminus A\) do not occur in this kernel's value. The sets
\(X\cap A,Y\) remain disjoint subsets of \(A\).

At branch 3 the kernel and its random set \(R\) are unchanged.
The set \(W\) lies in \(R\setminus X\) and excludes \(Y\),
because ancestry includes its targets. Thus the enlarged action
set is disjoint from the response. Its reference assignment exists
in the nonempty finite product \(S_W\). These vertices become
additional action arguments in the target, not conditioning variables
or prematurely fixed coordinates of the input kernel.

At branch 4 each child has the same valid kernel and graph, and
its response \(D_j\) is disjoint from action set \(R\setminus D_j\).
At branch 7, district recovery identifies the extracted kernel as
\(Q[D]\), with positive values and normalization by the support
and restriction lemmas. Its fixed set is \(N\setminus D\),
including the original values on \(R\setminus D\). The sets
\(X\cap D,Y\) are disjoint and belong to \(D\), because
\(Y\subseteq S\subseteq D\). This proves the invariant for
every child, without assuming the correctness of a returned answer.

It remains to justify that the listed branches exhaust all cases.
If a call reaches branch 3, then \(X\ne\varnothing\) and
\(R=\operatorname{an}_{G[R]}(Y)\). The response set cannot be
empty there: its ancestry would be empty, forcing \(R\) empty
and hence \(X\) empty. Thus \(Y\ne\varnothing\) and
\(S=R\setminus X\) is nonempty. Its district partition has
at least one member. Branch 4 handles every case with two or more;
otherwise its sole district is exactly \(S\).

A bidirected-connected set \(S\) in \(G[S]\) lies in exactly
one maximal district \(D\) of \(G[R]\), since its paths remain
paths in the larger graph. If the whole graph has one district,
branch 5 returns failure. Otherwise, equality \(D=S\) is handled
by branch 6. In the remaining case, \(S\subsetneq D\), and
\(D\subsetneq R\) because there is another maximal district.
This is the unique branch-7 district. All branches are therefore
defined on every reached call; empty \(R\) is handled at branch 1.

At branch 5, the failures of branches 1 and 2 give nonempty \(X\)
and \(R=\operatorname{an}_{G[R]}(Y)\). Failure of branch 3
gives \(S=\operatorname{an}_{G[S]}(Y)\). The one-district
tests give the two stated connectivity facts, and
\(\varnothing\ne Y\subseteq S\subsetneq R\) follows from
the preceding argument. The failure record stores these exact graphs
and sets, rather than asserting that their directed edges have already
been pruned into the source's confounded forests. Propagation retains
the local frame and all parent frames, so none of these facts or
their relationship to the original call is lost.
\end{proof}
\end{annexitem}

\begin{annexitem}{annex:p16.id-termination}{thm:p16-id-termination}{A strictly decreasing measure for every recursive edge}
\begin{proof}
At branch 2 the new random set \(A\) is a proper subset of
\(R\), and its nonaction set is
\(A\setminus(X\cap A)=A\setminus X\subseteq R\setminus X\).
Thus the first summand of \(\nu\) decreases and the second
does not increase. At branch 3 the random set stays \(R\),
while the nonaction set becomes
\((R\setminus X)\setminus W\), a proper subset because
\(W\ne\varnothing\).

At branch 4 a child retains \(R\), but its nonaction set is
\(D_j\), a proper subset of \(S=R\setminus X\) since at
least two nonempty districts partition \(S\). At branch 7 the
new random set \(D\) is a proper subset of \(R\), and
\[
 D\setminus(X\cap D)=D\setminus X=R\setminus X=S,
\]
because \(S\subseteq D\). Again one summand decreases and
the other is unchanged. These are all recursive branches; every
other branch returns directly.

The nonnegative integer measure therefore decreases strictly along
every call path. Its initial value is at most \(2|N|\), so a
path has at most \(2|N|+1\) calls. At each node there are
finitely many children and each finite graph, sum, product or
positive-denominator ratio operation is well defined. Induction on
\(\nu\), using finiteness of the children at a splitting node,
proves that the whole computation tree terminates. Failure propagation
can only shorten the executed part of this finite tree.
\end{proof}
\end{annexitem}

\begin{annexitem}{annex:p16.id-success-soundness}{thm:p16-id-success-soundness}{Soundness of each returning branch and its source derivation}
\begin{proof}
Induct on the proved measure \(\nu(R,X)\). The invariant identifies
the input kernel with the actual current model's law \(Q[R]\)
for every compatible original model. At branch 1 there is no
additional action, so ordinary marginalization is precisely the
target \(\mathcal T(R;Y,\varnothing)\).

At branch 2 the smaller call is valid and, if successful, its answer
is correct by induction. The second identity in
\cref{thm:p16-id-do-utility}, with its complete source Lemma 5
derivation, equates that target to the current one. The marginalized
kernel and the result are independent of the removed coordinates.

At branch 3 the child is smaller in measure and, if successful,
correctly computes the effect under \(X\cup W\). The third
utility identity, proved by the source Lemma 6 route, equates it
to the original effect for every \(w\). Substitution of the
recorded reference tuple therefore preserves the target and removes
any apparent need to supply an extra action value.

At branch 4 induction identifies each successful child as
\(Q[D_j]\): its response set is the whole nonaction set, so
there is no additional marginalization in that child's target.
The first utility identity, proved by the source Lemma 4 route,
states that the product of these kernels, summed over \(S\setminus Y\),
is the current target. The coordinate-summation convention makes
this a single consistent finite product and sum, even though a
coordinate was an action parameter in one child's kernel and is
summed in the parent expression. No independence of the observed
district tuples is asserted.

Branch 5 returns no probability expression, so it makes no success
claim. At branch 6 the fourth utility identity, the source Lemma 7
base, gives exactly the displayed marginal of the extracted kernel.
At branch 7 induction gives the target under \(X\cap D\) in
the model with kernel \(Q[D]\). The fifth utility identity,
the source Lemma 8 recursion, equates it to the current target at
the retained original action values on \(R\setminus D\).
This proves all success branches.

Every kernel operation is a finite marginal, product or conditional
ratio; all denominators are positive by the invariant. Substituting
the finite expressions from smaller calls therefore yields a finite
expression in the original observed law. For any two compatible
models sharing this law and its fixed value sets, the same graph
branches, kernels and expressions occur. A successful expression
equals each model's target, proving identification on this domain.

The utility theorem also gives the distinct do-calculus derivation:
each successful branch rewrites its target by the appropriate
source Lemma 4--8 sequence, while every computed district kernel
has the source's conditional-factor derivation. Substituting the
child derivations into the parent equations gives a finite chain
of ordinary probability operations, combination of disjoint fixed
actions, and the three sound rules, ending in the observed expression.
This follows by the same finite recursion induction. It is an
existence-of-derivation conclusion for successful runs, not the
still-separate assertion that all identifiable effects yield success.
\end{proof}
\end{annexitem}

\begin{annexitem}{annex:p16.id-worked-success}{thm:p16-id-worked-success}{The complete positive success traces}
\begin{proof}
\paragraph{The first model and every table entry.}
The sixteen latent tuples have product masses
\(\frac12\alpha(e_A)\beta(e_B)\gamma(e_C)>0\), summing to one.
For fixed observed \((a,b,c)\) and latent \(u\), the three
private inputs are forced to be
\(e_A=a\oplus u\), \(e_B=b\oplus a\), and
\(e_C=c\oplus\mathbf1_{\{b=1\text{ or }u=1\}}\).
Consequently
\[
 p(a,b,c)=\beta(b\oplus a)g(a,c;b),
\]
where
\[
 g(a,c;b)=\tfrac12\sum_{u=0}^1\alpha(a\oplus u)
                 \gamma(c\oplus\mathbf1_{\{b=1\text{ or }u=1\}}).
\]
At \(b=0\), the value of \(g\) is
\(\frac12[(4/5)(7/8)+(1/5)(1/8)]=29/80\) when \(a=c\),
and \(\frac12[(4/5)(1/8)+(1/5)(7/8)]=11/80\) otherwise.
At \(b=1\), summing \(\alpha(a\oplus u)\) gives one,
so \(g=1/16\) at \(c=0\) and \(g=7/16\) at \(c=1\).
Multiplying these four cases by \(\beta(b\oplus a)\) gives
the eight displayed numerators
\(87,33,5,35,11,29,15,105\) over 320. They sum to 320
and each is positive. Their prefix marginals are
\(p(a)=1/2\) and \(p(a,b)=\frac12\beta(b\oplus a)>0\).

\paragraph{Root call and the B factor.}
Start with \(R=\{A,B,C\}\), response \(\{C\}\), action
\(\{A\}\), and kernel \(p\), with no fixed outside coordinates.
Branch 1 does not apply. All three vertices are ancestors of \(C\),
so branch 2 does not apply. In \(S=\{B,C\}\), both vertices
are ancestors of \(C\), so branch 3 adds nothing. Its two
districts are \(\{B\},\{C\}\), so branch 4 calls for both factors.

The B child is \(\mathsf{ID}(\{B\},\{A,C\},p,G)\).
Branch 2 restricts to \(\{A,B\}\) with kernel
\(q_{AB}=\Sigma_Cp=\frac12\beta(b\oplus a)\), action \(\{A\}\),
and fixed coordinate \(C\), which the kernel ignores. In this
call branch 2 keeps both ancestors, branch 3 adds nothing, and
branch 4 has the single remaining district \(\{B\}\). The
current graph has two maximal districts \(\{A\},\{B\}\),
so branch 5 does not fail and branch 6 returns
\[
 \operatorname{Ext}^{AB}_{\{B\}}(q_{AB})
   =\frac{q_{AB}(a,b)}{\sum_{b'}q_{AB}(a,b')}
   =\beta(b\oplus a)=p(b\mid a).
\]
Every denominator in this child is \(1/2\).

\paragraph{The C factor and the larger-district recursion.}
The C child is \(\mathsf{ID}(\{C\},\{A,B\},p,G)\).
It retains all original ancestors at branch 2 and adds nothing
at branch 3. Its sole remaining district is \(\{C\}\).
The current maximal districts are \(\{A,C\}\) and \(\{B\}\),
so branches 5 and 6 do not return. Branch 7 extracts
\[
 q_{AC}=p(a)p(c\mid a,b)
       =\frac12\frac{p(a,b,c)}{p(a,b)}=g(a,c;b).
\]
This is a positive normalized kernel on \(\{A,C\}\), with
\(B=b\) retained as a fixed parameter. The new call has
response \(\{C\}\), action \(\{A\}\), and graph consisting
of \(A\leftrightarrow C\) with no directed edge.

Branch 2 now restricts to the sole directed ancestor \(\{C\}\).
It creates
\[
 q_C(c;b)=\sum_{a'}q_{AC}(a',c;b)
        =\sum_{a'}p(a')p(c\mid a',b),
\]
with fixed coordinates \(A,B\), of which \(A\) is ignored,
and with empty action set. Branch 1 returns this kernel. Its
response-one value is \((11+29)/80=1/2\) at \(b=0\), and
\(7/16+7/16=7/8\) at \(b=1\). Its response-zero values are
\(1/2\) and \(1/8\), so this last kernel is also positive.
The dummy \(a'\) is distinct from the original action argument.

The root combines the two successful children and sums over \(B\),
returning the displayed expression. At response one it gives
\[
 \tfrac34\tfrac12+\tfrac14\tfrac78=\tfrac{19}{32},\qquad
 \tfrac14\tfrac12+\tfrac34\tfrac78=\tfrac{25}{32}.
\]
Direct equation replacement confirms the values: after \(A=a\),
\(B=a\oplus E_B\) is independent of the fair \(U\). The
probability that \(B=1\) or \(U=1\) is \(5/8\) or \(7/8\).
Flipping that bit with independent probability \(1/8\) gives
\(1/8+(3/4)(5/8)=19/32\) or
\(1/8+(3/4)(7/8)=25/32\).

\paragraph{The second run and branch 3.}
In the second model, observationally \(A=U\) and
\(B=A\oplus U\oplus E_B=E_B\). The two variables are
independent fair bits; then \(C=B\oplus E_C\). Thus
\(p(a,b,c)=\frac14\eta(c\oplus b)>0\) as stated.
The observed graph has directed path \(A\to B\to C\) and
bidirected edge \(A\leftrightarrow B\), with total functions
using those declared inputs.

For response \(\{C\}\) and action \(\{B\}\), branch 2 keeps
all three directed ancestors. The nonaction graph on \(\{A,C\}\)
has only \(C\) as an ancestor of the response, so branch 3
sets \(W=\{A\}\), chooses \(a^\circ\), and calls with action
\(\{A,B\}\) and the same kernel. In that child, branch 2
again keeps the original ancestors, branch 3 has no additional
vertex, and the single remaining district is \(\{C\}\).
The full graph's maximal districts are \(\{A,B\}\), \(\{C\}\),
so branch 6 returns
\[
 \operatorname{Ext}^{ABC}_{\{C\}}(p)
       =p(c\mid a,b)=\frac{\frac14\eta(c\oplus b)}{1/4}
       =\eta(c\oplus b).
\]
The parent substitutes \(a^\circ\), which does not affect this
expression. Directly fixing \(B=b\) in the equations also gives
\(C=b\oplus E_C\), with this same law. This completes every
call and support calculation in the two successful runs.
\end{proof}
\end{annexitem}

\begin{annexitem}{annex:p16.id-worked-failure}{thm:p16-id-worked-failure}{The exact raw failure pair in the positive bow case}
\begin{proof}
The initial random set is \(R=\{A,Y\}\), the action is
\(\{A\}\), the response is \(\{Y\}\), and the fixed set
is empty. Its kernel is the uniform positive observed law. Branch 1
does not apply. The directed edge \(A\to Y\) makes both vertices
ancestors of \(Y\), so branch 2 does not apply. The nonaction
set is \(S=\{Y\}\), already its own response-ancestor set,
so branch 3 adds nothing. It has one district, and the full graph
also has one district because of \(A\leftrightarrow Y\).
Branch 5 therefore returns exactly
\(\mathsf{FAIL}(G[\{A,Y\}],G[\{Y\}])\).

The record retains local action \(\{A\}\), response \(\{Y\}\),
random set \(\{A,Y\}\), nonaction set \(\{Y\}\), empty
fixed set, binary value sets, the uniform kernel, and this branch-5
root frame. No conditional-ratio failure or reference assignment
occurred.

For this particular case, the full positive-law construction in
\cref{thm:p16-latent-foundation-limits} uses independent fair
latent roots \(U,T\), equation \(A=U\), and either
\(Y=\mathbf1_{\{T=1\text{ or }U\oplus A=1\}}\) or
\(Y=T\). Both observed laws are this same uniform table.
Their response-one intervention probabilities are \(3/4\) and
\(1/2\), so these two models prove nonidentification in this
binary positive-law case. This is an independently specified
witness for the worked run, not an inference of a general
completeness theorem from the bare failure tag.
\end{proof}
\end{annexitem}
